\documentclass[french]{article}
\usepackage{verbatim}
\usepackage{amsmath,amsfonts,amssymb}
\usepackage{pgfplots}
\usepackage[utf8]{inputenc}
\usepackage[french]{babel}
\usepackage{pgf,tikz,tkz-tab}
\usepackage{mathrsfs}
\usetikzlibrary{arrows}
\usepackage{pstricks}
\pgfplotsset{compat=newest}
\usepackage{geometry}
\usepackage{eurosym}
\geometry{hmargin=2cm,vmargin=1.5cm}


\begin{document}
\begin{center}
\section*{Devoir à la maison - L'algorithme de Babylone ou méthode babylonienne}
\end{center}
 
\textit{En mathématiques, la méthode de Héron ou méthode babylonienne est une méthode efficace d'extraction de racine carrée, c’est-à-dire de résolution de l'équation $x^2 = a$. Le but ici est d'étudier cette méthode, dans le cas particulier ou $a=2$ puis dans un cadre plus général par la suite.}\\

Voici le texte historique, extrait de son œuvre Métrica, tome 1. \\


\textit{« Puisque alors les 720 n'ont pas le côté exprimable, nous prendrons le côté avec une très petite différence ainsi. Puisque le carré le plus voisin de 720 est 729 et il a 27 comme côté, divise les 720 par le 27 : il en résulte 26 et deux tiers. Ajoute les 27 : il en résulte 53 et deux tiers. De ceux-ci la moitié : il en résulte 26 2' 3'. Le côté approché de 720 sera donc 26 2' 3'. En effet 26 2' 3' par eux-mêmes : il en résulte 720 36', de sorte que la différence est une 36e part d'unité. Et si nous voulons que la différence se produise par une part plus petite que le 36', au lieu de 729, nous placerons les 720 et 36' maintenant trouvés et, en faisant les mêmes choses, nous trouverons la différence qui en résulte inférieure, de beaucoup, au 36'. » Remarque: 2’ signifie 1/2} \\


\section*{Interprétation géométrique}
L'objectif est de déterminer la meilleure approximation de $\sqrt{720}$ c'est-à-dire la longueur d'un coté du carré
tel que l'aire de ce carré soit $720$. Pour cela on commence par choisir l'entier naturel qui au carré est le plus proche de 720. Ici $27^2=729$ donc 27 est l'entier naturel le plus proche de $\sqrt{720}$ par suite. On sait que le rectangle de dimension $\dfrac{720}{27} \times 27$ est d'aire 720. On veut cependant que cette figure se rapproche plus d'un carré que d'un rectangle ainsi on va faire la moyenne entre $\dfrac{720}{27}$ et $27$ soit $\dfrac{1449}{54}=26+\dfrac{5}{6}=26+ \dfrac{1}{2}+ \dfrac{1}{3} $. Ainsi on réitère le procédé mais non plus avec $27$ mais notre approximation $\dfrac{1449}{4}$ etc. 


\section*{Présentation de l'algorithme}
On considère le programme de calcul suivant que l'on appliquera à un nombre réel strictement positif $a$:

\begin{enumerate}
\item[] Choisir un nombre $a$
\item[] Calculer son inverse
\item[] Multiplier cet inverse par 2
\item[] Ajouter le résultat obtenu au nombre $a$
\item[] Diviser le résultat par 2 \\
\end{enumerate}

\begin{enumerate}
\item Appliquer cet algorithme à 10 nombres différents de votre choix.
\item Appliquer cet algorithme à 2, on note $R_1$ son résultat, donner son résultat sous forme fractionnaire.
\item Appliquer cet algorithme à $R_1$, on note $R_2$ son résultat, donner son résultat sous forme fractionnaire. 
\item Appliquer cet algorithme à $R_2$, on note $R_3$ son résultat, donner son résultat sous forme fractionnaire. 
\end{enumerate}

\section*{Application récursive de l'algorithme à l'aide de la calculatrice}

On souhaite appliquer ce programme de calcul plusieurs fois de suite à partir du nombre $a=2$.

\begin{enumerate}
\item Si j'applique cet algorithme à un nombre $x \in \mathbb{R}^*_+$, exprimer le résultat obtenu en fonction de $x$.
\item Entrer 2 sur la calculatrice puis EXE (ou ENTER).
\item Faites la manipulation : $\dfrac{ANS + 2 \times \dfrac{1}{ANS}}{2}$ puis EXE(ou ENTER).
\item Faire cette manipulation 6 fois de suite et donner la liste des six résultats obtenus.
\end{enumerate}

\section*{Interprétation mathématique et comparaison}
Soit $a$ un nombre réel différent de $\sqrt2$.

\begin{enumerate}
\item Comparer les résultats de la question précédente à $\sqrt{2}$. Quelle(s) constatation(s) peut-on faire ?
\item Développer $(a-\sqrt{2})^2$.
\item En déduire que $a^2+2 > 2a\sqrt{2}$.
\item Justifier que pour $a>0$ , $\dfrac{a^2+2}{2a}> \sqrt{2}$.
\item En déduire que pour tout $a>0$ ; $\dfrac{1}{2} (a+ \dfrac{2}{a})>\sqrt{2}$.
\end{enumerate}

\section*{Une équation}

Résoudre l'équation d'inconnue $a$ strictement positif suivante : $$\dfrac{1}{2} (a+ \dfrac{2}{a})=a$$

\section*{Première scientifique}

Soit $U_{n+1}=f(U_n)$ avec $U_0=2$ où $f(x)=\dfrac{1}{2}(x+\dfrac{2}{x})$.

\begin{enumerate}
\item Quel est l'ensemble de définition, de continuité et de dérivabilité de la fonction $f$ ?
\item Etudier les variations de la fonction $f$.
\item Montrer par récurrence que $(U_n)_{n\in \mathbb{N}}$ est décroissante.
\item En déduire que $(U_n)_{n\in \mathbb{N}} $ converge.
\item Faire un programme sur Python qui donne les 50 premières itérations de $(U_n)_{n\in \mathbb{N}}$ en fonction de $U_0 \in \mathbb{R}^*_+ $
\end{enumerate}

\section*{Terminale scientifique}
On suppose que $(U_n)_{n\in \mathbb{N}} $ converge vers une limite notée $l$.
\begin{enumerate}
\item En passant à la limite l'équation $U_{n+1}=f(U_n)$  et en justifiant pourquoi, donner l'équation que respecte nécessairement $l$.
\item Résoudre l'équation de point fixe $f(x)=x$, en déduire en justifiant la limite de $(U_n)_{n\in \mathbb{N}} $.
\end{enumerate}

\section*{Terminale scientifique (Niveau I)}
On suppose vrai le résultat précédent et on pose ici $U_0=c$ où $c\in \mathbb{R}^*_+$.
\begin{enumerate}
\item Montrer que le résultat reste valable si $c\geq \sqrt{2}$.
\item En déduire de la question 4 de la partie "Interprétation mathématique et comparaison" que le résultat reste valable si $c>0$.
\end{enumerate}

\section*{Terminale scientifique (Niveau II)}

Soit $a\in \mathbb{R}_+^*$ $U_{n+1}=f_a(U_n)$ avec $U_0=a$ où $f_a(x)=\dfrac{1}{2}(x+\dfrac{a}{x})$

\begin{enumerate}
\item Quel est l'ensemble de définition, de continuité et de dérivabilité de la fonction $f_a$ ?
\item Étudier les variations de la fonction $f_a$.
\item Montrer par récurrence que $(U_n)_{n\in \mathbb{N}} $ est décroissante.
\item En déduire que $(U_n)_{n\in \mathbb{N}} $ converge.
On suppose que $(U_n)_{n\in \mathbb{N}} $ converge vers un limite notée $l_a$.
\item En passant à la limite l'équation $U_{n+1}=f_a(U_n)$  et en justifiant pourquoi, donner l'équation que respecte nécéssairement $l_a$.
\item Résoudre l'équation de point fixe $f_a(x)=x$, en déduire en justifiant la limite de $(U_n)_{n\in \mathbb{N}} $.\\

On suppose le résultat précédent et on pose ici $U_0=c$ où $c\in \mathbb{R}^*_+$.
\item Montrer que le résultat reste valable si $c\geq \sqrt{a}$.
\item Développer $(x-\sqrt{a})^2$ et déduire que $f_a(x)\geq \sqrt{a}$.
\item En déduire que le résultat reste valable si $c>0$.
\end{enumerate}

\section*{Terminale scientifique (Niveau III)}
\begin{enumerate}
\item Montrer que $U_{n+1}-\sqrt{a}=\dfrac{(U_n-\sqrt{a})^2}{2U_n}$
\item En déduire alors l'inégalité $\vert U_{n+1}-\sqrt{a} \vert \leq \dfrac{\vert U_n-\sqrt{a}\vert ^2}{2U_n}$ 
\item On suppose que $U_n \geq \sqrt{a}$  en déduire alors que $ \dfrac{\vert U_{n+1}-\sqrt{a} \vert}{\vert U_n-\sqrt{a}\vert ^2} \leq \dfrac{1}{2\sqrt{a}}$
\end{enumerate}

On dit que la suite $(U_n)_{n\in \mathbb{N}} $ est convergente d'ordre 2, on dit aussi que sa convergence est quadratique. Autrement dit, la précision de l'approximation double à chaque étape. On constate bien la convergence quadratique si $U_0=1$ et $a=2$ (2 décimales exactes au deuxième calcul, 5 au troisième, 11 au quatrième, 23 au cinquièmes...) 

\newpage 
\begin{center}
\section*{Élément de correction}
\end{center}
\renewcommand{\arraystretch}{2.5}

\subsection*{Présentation de l'algorithme}
\begin{enumerate}
\item 
\begin{center}
\begin{tabular}{|*{11}{c|}}
    \hline
    $a$ & 1  & 2  & 3  & 4  & 5  & 6  & 7  & 8  & 9  & 10 \\
    \hline
    $\dfrac{1}{a}$ & $\dfrac{1}{1}$  & $\dfrac{1}{2}$  & $\dfrac{1}{3}  $& $\dfrac{1}{4}  $& $\dfrac{1}{5} $& $\dfrac{1}{6}  $& $\dfrac{1}{7}  $& $\dfrac{1}{8}  $& $\dfrac{1}{9}  $&$ \dfrac{1}{10} $\\
    \hline

 $\dfrac{2}{a}$ & $\dfrac{2}{1}$  & $\dfrac{2}{2}$  & $\dfrac{2}{3}  $& $\dfrac{2}{4}  $& $\dfrac{2}{5} $& $\dfrac{2}{6}  $& $\dfrac{2}{7}  $& $\dfrac{2}{8}  $& $\dfrac{2}{9}  $&$ \dfrac{2}{10} $\\

    \hline
   $\dfrac{2}{a}+a$ & $\dfrac{2}{1}+1$  & $\dfrac{2}{2}+2$  & $\dfrac{2}{3} +3 $& $\dfrac{2}{4} +4 $& $\dfrac{2}{5} +5$& $\dfrac{2}{6} +6 $& $\dfrac{2}{7} +7 $& $\dfrac{2}{8} +8 $& $\dfrac{2}{9} +9  $&$ \dfrac{2}{10} +10 $\\
    \hline
    $\dfrac{2+a}{2a^2}$ & $\dfrac{3}{2}$  & $\dfrac{3}{2}$  & $\dfrac{11}{6} $& $\dfrac{5}{4} $& $\dfrac{27}{10}$& $\dfrac{19}{6} $& $\dfrac{51}{14}$& $\dfrac{33}{8} $& $\dfrac{83}{18} $&$ \dfrac{51}{10} $\\
    \hline
\end{tabular}
\end{center}
\item
\begin{center}
\begin{tabular}{|*{4}{c|}}
    \hline
    $a$ & $2$  & $\dfrac{3}{2}$  & $\dfrac{17}{12}$ \\
    \hline
    $\dfrac{1}{a}$  & $\dfrac{1}{2}$  & $\dfrac{2}{3}  $& $\dfrac{12}{17}  $\\
    \hline
 $\dfrac{2}{a}$ & $\dfrac{2}{2}$  & $\dfrac{4}{3}  $& $\dfrac{24}{17}  $\\
    \hline
   $\dfrac{2}{a}+a$  & $\dfrac{2}{2}+2$  & $\dfrac{17}{6} +3 $& $\dfrac{577}{204}$\\
    \hline
    $\dfrac{2+a}{2a^2}$ &  $R_1=\dfrac{3}{2}$  & $R_2=\dfrac{17}{12} $& $R_3=\dfrac{577}{408} $\\
    \hline
\end{tabular}
\end{center}

\end{enumerate}

\section*{Application récursive de l'algorithme à l'aide de la calculatrice}

\begin{enumerate}
\item Le résultat obtenu serai $\dfrac{2+x}{2x^2}$.
\item On trouve exactement $R_1,R_2$ et $R_3$ pour les 3 premières itérations qui sont respectivement égales à $1,5$, $\approx 1,41666666$ et $ \approx 1,41421569$ puis ensuite on reste à $\approx 1,41421356$ pour les 3 itérations suivantes.
\end{enumerate}

\section*{Interprétation mathématique et comparaison}
\begin{enumerate}

\item Quand on compare $\sqrt2$ à la dernière itération on trouve $-3.73095 \times 10^{-10}$. On remarque que la différence est proche de 0.  
\item $(a-\sqrt2)^2=a^2-2\sqrt2 a+ 2$
\item Comme $a\neq \sqrt2$, on a $(a-\sqrt2)^2 > 0$, d'où $a^2 -2a\sqrt2 +2 > 0$, ce qui est équivalent en ajoutant $2a\sqrt2$ de chaque coté de l'inégalité a $a^2 +2 > 2a\sqrt2$.
\item Comme $a>0$ on a $2a\neq0$ on divise alors l'inégalité par $2a$, l'ordre de l'inégalité est conservé car $2a$ est strictement positif et on obtient $\dfrac{a^2 +2}{2a}>\sqrt2$.
\item Comme $\dfrac{a^2+2}{2a}=\dfrac{1}{2}(\dfrac{a^2+2}{a})=\dfrac{1}{2}(\dfrac{a^2}{a}+\dfrac{2}{a})= \dfrac{1}{2}(a+\dfrac{2}{a})$ on a par la question précédente que $\dfrac{1}{2}(a+\dfrac{2}{a}) >\sqrt2$.
\end{enumerate}

\newpage

\section*{Une équation}
\begin{enumerate}
\item Soit $a>0$ tel que $$ \dfrac{1}{2}(a+\dfrac{2}{a})=a$$
$$\Longleftrightarrow a+\dfrac{2}{a}=2a $$
$$\Longleftrightarrow \dfrac{2}{a}=2 $$
$$\Longleftrightarrow 2=a^2  \; \; \textrm{car} \; a \neq 0$$
$$\Longleftrightarrow a=\sqrt2 \; \; \textrm{car} \; a>0$$

\end{enumerate}

\end{document}


